\chapter{Frontend và một host duy nhất}
\label{ch:frontend}

Cột mốc~7 đưa thứ có thể triển khai cuối cùng --- frontend React ---
vào cụm và giành lấy đường dẫn \code{/} trên \code{ts.local} mà
Ingress của Chương~18 đã cố ý để trống. Đây là một cột mốc ngắn tính
theo số dòng YAML và dài tính theo khái niệm: đó là image duy nhất mà
pipeline build từ Dockerfile, pod duy nhất không phải JVM, và là thay
đổi biến ``các service gọi được bằng \code{curl}'' thành ``một ứng
dụng dùng được trong trình duyệt''.

\section{Frontend là gì, dưới mắt người vận hành}

Một ứng dụng React đã build không phải là một server. \code{npm run
build} (kiểm tra TypeScript, rồi đóng gói bằng Vite) sinh ra một thư
mục file tĩnh: một \code{index.html}, một bundle JavaScript có mã băm
nội dung trong tên (\code{index-CqZING8N.js}), một file CSS, một
favicon. Không có gì trong đó chạy trên server; từng dòng của nó chạy
trong trình duyệt của khách, rồi trình duyệt gọi \code{/api/...} qua
HTTP như bất kỳ client nào khác. Vậy ``triển khai frontend'' nghĩa là
triển khai một \emph{file server} có chứa những file đó --- ở đây là
nginx --- và bảo router biên biết request nào là file, request nào là
API.

Cách nhìn đó quyết định mọi thứ còn lại trong chương này. Không config
server (ứng dụng không có thuộc tính nào để import), không Eureka
(không ai gọi nó theo tên), không secret (trình duyệt đọc được mọi
thứ frontend nắm giữ, nên không được có gì bí mật ở đó), và một dấu
chân tài nguyên nhỏ hơn cả chi phí nạp class của JVM nhỏ nhất.

\section{Hai giai đoạn, một Dockerfile}

Jib (Chương~13) build image JVM từ \code{target/classes}; nó không
chạy được \code{npm}. Vì thế frontend dùng Dockerfile nhiều giai đoạn
(multi-stage) kinh điển --- và lý do nó có hai giai đoạn cũng chính là
lý do các image JVM dùng image gốc JRE thay vì JDK:

\begin{lstlisting}[language=dockerfile]
FROM node:24-alpine AS build            (*@\co{1}@*)
WORKDIR /app
COPY package.json package-lock.json ./  (*@\co{2}@*)
RUN npm ci --no-audit --no-fund
COPY . .
RUN npm run build

FROM nginx:1.27-alpine                  (*@\co{3}@*)
COPY nginx.conf /etc/nginx/conf.d/default.conf
COPY --from=build /app/dist /usr/share/nginx/html   (*@\co{4}@*)
EXPOSE 80
\end{lstlisting}

\begin{description}
  \item[\co{1}] Giai đoạn build có Node, trình biên dịch TypeScript, và
  sau \code{npm ci} là chừng 300\,MB \code{node\_modules}. Không thứ
  nào trong đó cần cho việc \emph{phục vụ} kết quả.
  \item[\co{2}] Lockfile trước, mã nguồn sau --- quy tắc cache theo
  layer của Chương~12. \code{npm ci} cài đúng các phiên bản đã ghim
  trong \code{package-lock.json} và từ chối sửa nó; một chỉnh sửa mã
  nguồn chỉ làm mất hiệu lực các layer từ \code{COPY . .} trở xuống,
  nên việc tải dependency được bỏ qua ở mọi lần rebuild không đụng
  tới \code{package.json}.
  \item[\co{3}] Một image gốc mới tinh. Image cuối cùng là nginx cộng
  file tĩnh: khoảng 15\,MB, không Node, không trình biên dịch, không
  trình quản lý gói --- bề mặt tấn công nhỏ hơn và kéo nhanh hơn, miễn
  phí.
  \item[\co{4}] \code{--from=build} là phép sao chép multi-stage: chỉ
  \code{dist/} băng từ giai đoạn một sang giai đoạn hai. Mọi thứ khác
  ở giai đoạn một bị vứt bỏ khi build kết thúc.
\end{description}

\begin{decision}[Sao không thêm Node vào image Jenkins rồi dùng Jib?]
Jib có thể gói bất kỳ thư mục nào thành image, nên con đường ``Node
trong Jenkins, Jib để đóng gói'' là có. Nó bị loại vì che giấu kỹ
thuật build image phổ biến nhất sau kỹ thuật ít phổ biến nhất. Mọi
đội ngũ phát hành frontend đều có Dockerfile này; đọc và viết được nó
tự thân đã là một kỹ năng thị trường tuyển dụng cần, và có một image
không-Jib trong pipeline cũng chứng minh pipeline không âm thầm mang
hình dáng của Jib. Cái giá là một Docker CLI trong image Jenkins và
một socket daemon được mount --- Mục~\ref{sec:dockersock}.
\end{decision}

\section{nginx: fallback giúp SPA hoạt động}

\begin{lstlisting}[language=bash]
server {
    listen 80;
    root  /usr/share/nginx/html;
    location / {
        try_files $uri $uri/ /index.html;          (*@\co{1}@*)
    }
    location /assets/ {
        add_header Cache-Control "public, max-age=31536000, immutable"; (*@\co{2}@*)
    }
    location = /healthz { return 200 "ok\n"; }      (*@\co{3}@*)
}
\end{lstlisting}

\begin{description}
  \item[\co{1}] Dòng mà mọi ứng dụng một trang (SPA) đều cần và mọi lần
  triển khai đầu tiên đều quên. React Router (\code{BrowserRouter}
  trong \code{main.tsx}) sở hữu những URL như \code{/courses/42} hoàn
  toàn trong trình duyệt: bấm một link không bao giờ hỏi server. Nhưng
  một lần refresh, một bookmark hay một link dán vào \emph{có} hỏi
  server về \code{/courses/42}, file đó không tồn tại, và một file
  server thuần trả 404. \code{try\_files} thử đường dẫn nguyên văn, rồi
  thư mục, rồi rơi về \code{index.html} --- ứng dụng khởi động và
  router của nó đọc URL. OAuth callback (\code{/oauth/callback}) cũng
  dựa vào điều này: redirect của Google là một request lạnh tới một
  route chỉ tồn tại trong JavaScript.
  \item[\co{2}] Vite ghi mã băm nội dung vào tên mọi asset, nên file
  đổi thì luôn có tên mới. Điều đó khiến ``cache vĩnh viễn'' an toàn
  cho \code{/assets/} --- và \emph{không an toàn} cho \code{index.html},
  thứ phải được tải lại để biết các tên mới sau mỗi lần deploy. Cache
  file HTML là cách người dùng rơi vào cảnh ứng dụng cũ đi đòi những
  bundle mà server không còn nữa.
  \item[\co{3}] Một đích cho probe không tốn gì và không bao giờ chạm
  đĩa (các probe của Chương~17 cần thứ gì đó để gọi).
\end{description}

Hãy để ý điều file này \emph{không} làm: proxy \code{/api}. Trên
laptop, \code{vite.config.ts} proxy \code{/api} tới gateway ở
\code{localhost:8080} để dev server và API trông như một origin. Trong
cụm, việc đó thuộc về Traefik --- nginx không bao giờ thấy một request
API nào cả.

\section{Giành lấy \code{/}: một origin}

\begin{lstlisting}[language=yaml]
kind: Ingress
metadata: { name: frontend }
spec:
  rules:
    - host: ts.local
      http:
        paths:
          - path: /                     (*@\co{1}@*)
            pathType: Prefix
            backend: { service: { name: frontend, port: { number: 80 } } }
\end{lstlisting}

\begin{description}
  \item[\co{1}] Một quy tắc bắt-tất-cả, trong một object Ingress
  \emph{thứ hai} đặt cạnh Ingress của gateway. Traefik gộp mọi quy
  tắc Ingress của một host và định tuyến theo \emph{tiền tố khớp dài
  nhất}: \code{/api/courses} khớp cả \code{/} lẫn \code{/api}, và
  \code{/api} thắng. Ba tiền tố của gateway vì thế tiếp tục hoạt động
  không đổi, còn mọi thứ khác --- \code{/}, \code{/login},
  \code{/courses/42}, \code{/assets/...} --- rơi xuống nginx. Hai
  object, một bảng định tuyến gộp; quy tắc ``apply không bao giờ xóa''
  của Chương~20 vẫn áp dụng cho từng cái.
\end{description}

Hệ quả chính là lý do cột mốc này tồn tại: \textbf{một origin}. Trình
duyệt tải trang từ \code{http://ts.local} và gọi
\code{http://ts.local/api/...} --- cùng scheme, host và cổng. So với
Compose, nơi trang đến từ \code{localhost:3000} còn API từ
\code{localhost:8080}: khác cổng là khác origin, nên mọi lời gọi API
đều có một CORS preflight đi trước và \code{CorsConfig} của gateway
phải liệt kê rõ \code{localhost:3000}. Trong cụm, \code{CorsConfig}
đơn giản là không bao giờ được hỏi tới --- chẳng có request
cross-origin nào để kiểm soát. Cookie, redirect của OAuth, và bất kỳ
chính sách \code{SameSite} nào về sau đều đơn giản đi vì cùng lý do.

\begin{figure}[h]
\centering
\begin{tikzpicture}[node distance=5mm and 8mm]
  \node[boxdark] (b) {Trình duyệt};
  \node[box, right=22mm of b] (t) {Traefik};
  \node[box, right=34mm of t, yshift=9mm] (fe) {frontend\\\scriptsize nginx};
  \node[box, right=34mm of t, yshift=-9mm] (gw) {api-gateway};
  \draw[flow] (b) -- node[lbl,above]{\scriptsize Host: ts.local} (t);
  \draw[flow] (t.east) -- node[lbl,above,sloped,pos=0.5]{\scriptsize /} (fe.west);
  \draw[flow] (t.east) -- node[lbl,below,sloped,pos=0.5]{\scriptsize /api, /oauth2, /login/oauth2} (gw.west);
\end{tikzpicture}
\caption{Tiền tố dài nhất thắng: một host, hai backend.}
\end{figure}

Một thuộc tính đi theo việc đổi host. \code{OAuth2LoginSuccessHandler}
kết thúc đăng nhập Google bằng cách chuyển hướng trình duyệt tới
\code{app.oauth2.success-redirect-uri}, mà \code{auth-service.yml}
của config-server đặt là \code{localhost:3000/oauth/callback} --- đúng
trên laptop, sai trong cụm. \code{auth-service.yaml} ghi đè nó bằng
biến môi trường
\code{APP\_OAUTH2\_SUCCESS\_REDIRECT\_URI=http://ts.local/oauth/callback}
--- quy tắc ưu tiên của Chương~16 (biến môi trường thắng config
server) làm đúng việc nó đã làm cho URL datasource.

\section{Manifest}

Deployment mang hình dáng của Chương~15 với gần như mọi thứ bị lược
bỏ: không \code{envFrom}, không biến \code{SPRING\_*}, không
\code{startupProbe} (nginx lên trong vài mili giây), và
\code{requests: \{ memory: 32Mi, cpu: 10m \}} --- một phần mười hai
của JVM nhỏ nhất. \code{imagePullPolicy: Always} với tag
\code{:latest} làm mặc định của manifest, được pipeline ghi đè bằng
tag commit y hệt các service JVM.

\section{Pipeline mọc thêm một stage}
\label{sec:dockersock}

\begin{lstlisting}
stage('Frontend image') {
  steps {
    dir('Project/TeacherSupporter') {
      sh '''
        docker build -t ${REGISTRY}/frontend:${TAG} \
                     -t ${REGISTRY}/frontend:latest frontend   (*@\co{1}@*)
        docker push ${REGISTRY}/frontend:${TAG}
        docker push ${REGISTRY}/frontend:latest
      '''
    }
  }
}
\end{lstlisting}

\begin{description}
  \item[\co{1}] Build context là \code{frontend/}, được cắt gọn bởi
  một \code{.dockerignore} loại trừ \code{node\_modules} và
  \code{dist}: các module native của laptop Windows không được gửi
  lên daemon, và image tự build lấy của mình. \code{frontend} gia nhập
  \code{SERVICES} (danh sách của kubectl) nhưng không vào
  \code{MODULES} (của Maven) --- các vòng lặp \code{set image} /
  \code{rollout status} của stage Deploy xử lý nó mà không cần đổi gì.
\end{description}

Để \code{docker build} chạy được bên trong container Jenkins, hai thứ
thay đổi trong thiết lập của Chương~22. Image có thêm Docker
\emph{CLI} (\code{docker-ce-cli} --- chỉ phần client), và
\code{run.sh} mount \code{/var/run/docker.sock} của host vào
container. CLI nói chuyện qua socket đó với daemon \emph{của host};
không có daemon nào bên trong Jenkins và không có Docker-in-Docker.
Image tạo ra nằm trong kho image của host, và \code{docker push
localhost:5000/...} tới được registry qua mạng host mà container vốn
đã dùng chung.

\begin{pitfall}[permission denied trên /var/run/docker.sock]
Triệu chứng: CLI đã cài, socket đã mount, và mọi lệnh \code{docker}
đều thất bại với lỗi quyền. Nguyên nhân: socket thuộc sở hữu của nhóm
\code{docker} trên host --- GID 984 trên \code{ts-server} --- và con số
đó chẳng có ý nghĩa gì bên trong image, nơi user \code{jenkins} thuộc
một nhóm cùng tên nhưng GID khác (hoặc không có). Cách sửa trong
\code{run.sh}: \code{--group-add "\$(stat -c \%g /var/run/docker.sock)"}
thêm GID số của host vào user của container lúc chạy, không phải nướng
một con số đặc thù của máy vào Dockerfile.
\end{pitfall}

\begin{decision}[Socket là root]
Mount socket Docker trao quyền kiểm soát host tương đương root: ai
chạy được \code{docker run -v /:/host} là sở hữu cả máy. Với một
homelab một người dùng mà Jenkins chỉ truy cập được qua tailnet, đây
là sự đánh đổi được chấp nhận, nói rõ thành lời. Production build
image mà không cần daemon nào cả --- Kaniko hoặc BuildKit chạy như
một pod không đặc quyền, hoặc Jib cho JVM --- và việc chuyển Jenkins
vào cụm ở Chương~27 là nơi khoản thỏa hiệp này sẽ được trả lại.
\end{decision}

\begin{hardware}
Pod frontend rẻ nhất namespace (yêu cầu 32\,Mi), nhưng \emph{build}
thì không: \code{npm ci} cộng \code{tsc} cộng Vite chạy gần một
CPU-phút và một gigabyte RAM trên controller, trong khi Maven có thể
vẫn đang chạy trong cùng container 2\,GB đó. Jenkinsfile giữ các stage
tuần tự vì lý do đó; trên cái máy này, stage song song chỉ tranh giành
nhau.
\end{hardware}

\section{Kiểm chứng cột mốc}

\begin{handson}[danh sách kiểm tra cột mốc 7]
\begin{enumerate}
  \item \code{kubectl -n ts get ingress} --- hai object trên
  \code{ts.local}; \code{kubectl -n ts get pods -l app=frontend} Ready.
  \item Mở \code{http://ts.local/} --- trang đăng nhập. Rồi dán thẳng
  \code{http://ts.local/courses} vào thanh địa chỉ: một trang, không
  phải 404 của nginx --- fallback SPA hoạt động.
  \item Đăng ký, đăng nhập, tạo một khóa học: tab network của trình
  duyệt cho thấy các request \code{/api/...} không có preflight
  \code{OPTIONS} nào --- cùng origin.
  \item \code{curl -sI http://ts.local/assets/index-*.js | grep
  Cache-Control} --- \code{immutable}; làm tương tự với \code{/} ---
  không có header cache.
  \item Đổi một đoạn chữ bất kỳ trong một trang, push, và xem pipeline:
  stage \code{Frontend image} dùng lại layer \code{npm ci}, và rollout
  thay vào một image 15\,MB.
\end{enumerate}
\end{handson}

Build~\#9 là lần chạy đầu tiên của stage mới: \code{docker build}
hoàn tất trong 90 giây trên controller, rollout thay
\code{frontend:f5d9213} vào, và danh sách kiểm tra ở trên đạt đúng
như đã viết --- \code{/courses} trả về \code{index.html},
\code{/api/courses} trả về JSON 401 của gateway, asset mang
\code{immutable} và \code{/} không mang header cache nào.

\begin{pitfall}[Ingress sống lâu hơn manifest của nó]
\code{kubectl get ingress} hôm đó liệt kê \emph{ba} object trên
\code{ts.local}: của gateway, của frontend --- và
\code{course-service}, tám ngày tuổi, định tuyến \code{/courses} thẳng
tới service mà không kiểm tra JWT. Manifest của nó đã bị gỡ khỏi
\code{course-service.yaml} từ Chương~18 kèm một comment nói rõ điều
đó, nhưng quy tắc của Chương~20 vẫn đúng: \code{apply} không bao giờ
xóa thứ nó không còn thấy. Phải cần một lệnh tay
\code{kubectl -n ts delete ingress course-service}. Bài học không phải
là câu lệnh mà là việc kiểm toán: sau mỗi cột mốc, hãy liệt kê thứ
cụm \emph{có}, không chỉ thứ git \emph{nói} --- cho đến khi việc dọn
(prune) của GitOps (Chương~27) biến hai thứ đó thành một.
\end{pitfall}

\section{Truy cập từ bất cứ đâu: tên trên tailnet}

Mọi thứ ở trên trả lời trên \code{ts.local} --- một cái tên chỉ tồn
tại trong file hosts của một chiếc laptop, trên một LAN. ``Live''
nghĩa là một URL phân giải được trên mọi thiết bị bạn có, và cái máy
đã có sẵn một cái: Tailscale cho mỗi node một tên MagicDNS,
\code{ts-server.tailbfe002.ts.net}, mà mọi peer trên tailnet phân
giải được mà không cần sửa file hosts. Traefik định tuyến theo header
\code{Host} (Chương~18), nên một cái tên nó không liệt kê sẽ nhận 404
dù gói tin có tới nơi. Cách sửa là thêm một mục \code{rules} nữa với
cùng danh sách đường dẫn trong cả hai object Ingress (Phụ lục~D có
chúng) --- apply bằng tay với một bộ lọc chỉ trích ra các tài liệu
Ingress từ đầu ra của \code{kubectl kustomize}, vì một lệnh
\code{apply -k} đầy đủ sẽ đặt lại tag image của mọi Deployment về
\code{1.0.0-SNAPSHOT} của manifest và roll cả namespace (\code{set
image} của Chương~23 nằm \emph{sau} \code{apply} chính vì lý do này).

Cái tên đó rồi phơi ra một lỗi mà LAN đã che giấu. Một lần đăng ký
qua host mới sinh ra mail kích hoạt (chuỗi của Chương~25, trọn vẹn
qua tailnet), nhưng link bên trong ghi
\code{http://localhost:8080/api/auth/activate?code=...}: một host
không điện thoại nào tới được, và một endpoint \code{POST} mà
\code{GET} của trình đọc mail đằng nào cũng không dùng nổi.
\code{EmailService} dựng mọi link --- kích hoạt, mời, đăng nhập --- từ
chuỗi literal. Sửa chữa là quy tắc của Chương~16 áp dụng cho URL:
origin công khai là \emph{cấu hình}. \code{app.frontend.base-url} mặc
định là dev server của Vite trong \code{notification-service.yml} của
config-server; manifest ghi đè nó bằng \code{APP\_FRONTEND\_BASE\_URL}
đặt thành tên tailnet, cùng thủ thuật relaxed binding như
\code{SPRING\_MAIL\_HOST}; và các link giờ trỏ tới trang
\code{/activate} của frontend, nơi thực hiện \code{POST}. Hai biến
OAuth chuyển sang cùng host, đồng nghĩa với thêm một mục nữa dưới
``Authorized redirect URIs'' của Google --- việc đăng ký mà Chương~5
gọi là thứ duy nhất không manifest nào tự động hóa được.

\begin{pitfall}[một link là một hợp đồng với bên ngoài]
Bất cứ thứ gì rời khỏi hệ thống --- một email, một payload webhook,
một mã QR --- đều mang một URL phải hoạt động từ \emph{bên ngoài}, nên
nó không bao giờ được là tên Service và không bao giờ nên là literal.
Bài kiểm tra rất rẻ: đăng ký từ một thiết bị không phải laptop build
và bấm vào. Trên stack này, chính request đầu tiên qua hostname mới
đã phơi ra một lỗi có mặt từ commit đầu tiên.
\end{pitfall}

\section{Ngoài dự án}

Mẫu hình --- router biên tách tĩnh khỏi API theo đường dẫn, phần tĩnh
đến từ một file server ngờ nghệch --- là phổ quát, với file server
thường được thay bằng lưu trữ đối tượng cộng CDN (S3 + CloudFront,
hoặc một static host như Netlify) và việc tách do quy tắc origin của
CDN làm. Giai đoạn build vẫn y nguyên nửa đầu của Dockerfile này. Hai
thói quen chuyển sang trực tiếp: băm-và-cache-vĩnh-viễn cho asset kèm
\code{index.html} không cache, và fallback SPA --- theo cách nói của
CloudFront là ``ánh xạ 403/404 sang \code{/index.html} với mã 200'',
cùng ý tưởng nhưng dùng vụng hơn.

\begin{quiz}
\begin{enumerate}
  \item Vì sao image frontend có hai dòng \code{FROM}, và chính xác
  thứ gì băng qua giữa chúng?
  \item Một người dùng bookmark \code{http://ts.local/courses/7} và mở
  nó sáng hôm sau. Lần theo request qua Traefik và nginx, và nói dòng
  nào khiến nó thành công.
  \item Cả hai object Ingress đều khai báo \code{host: ts.local}.
  Traefik quyết định thế nào để \code{/api/courses} đi tới gateway và
  \code{/assets/x.js} đi tới nginx?
  \item Giải thích vì sao \code{CorsConfig} của gateway trở nên không
  liên quan trong cụm trong khi vẫn cần thiết trên laptop.
  \item \code{--group-add "\$(stat -c \%g /var/run/docker.sock)"} sửa
  điều gì, và vì sao giá trị được tính ra thay vì viết sẵn?
  \item Link trong mail kích hoạt đã sai hàng tháng mà không ai để ý.
  Nêu hai khiếm khuyết riêng biệt trong
  \code{http://localhost:8080/api/auth/activate?code=...}, và nói mỗi
  cái sẽ hỏng ở môi trường nào (laptop, compose, cụm).
\end{enumerate}
\end{quiz}
